\SourcePage{384}
\SGAchapter{X}{平展上同调中的 Euler--Poincaré 公式}{Formule d'Euler--Poincaré en cohomologie étale}
\begin{center}A. Grothendieck\\（I. Bucur 整理）\end{center}
本讲的目的是陈述并证明 Euler--Poincaré 型的公式（7.2）。它推广了 Ogg 和 Chafarevitch 各自独立得到的公式；M. Raynaud 的报告 [1] 从\Term{etale-cohomology}{平展上同调}的角度介绍了他们的结果。

\Term{euler-poincare-formula}{Euler--Poincaré 公式}实质上使用了\Term{derived-category}{导出范畴}的\Term{form}{形式}体系。与上述文献一样，证明要经过 Weil 的一个著名公式的变形（5.1）。

\SGAsection{1}{带作用的概形上的层}
设 $X$ 是一个\Term{scheme}{概形}，群 $G$ 通过\Term{automorphism}{自同构}\Term{action}{作用}于其上。仿照（[T] 第 V 章），可以定义 $X$ 上以 $G$ 为\Term{operator-group}{算子群}的\Term{etale-sheaf}{平展层}，简称“$(X_{\mathrm{ét}},G)$ 上的\Term{sheaf}{层}”。

按照定义，给出这样的\Term{etale-sheaf}{平展层}，就是给出一个\Term{sheaf}{层} $F\in\operatorname{Ob}(\widetilde X_{\mathrm{ét}})$，并对每个 $g\in G$ 给出一个\Term{morphism}{态射}
\[
S_g:F\longrightarrow g^*(F),
\]
\SourcePage{385}
使下列条件成立：
\[
\text{a)}\quad S_e=1_F;
\]
\noindent b) 对 $g,h\in G$，下图交换：
\[
\begin{tikzcd}
F\arrow[r,"S_{hg}"]\arrow[d,"S_h"']&(hg)^*(F)\arrow[d,"{C_{h,g}(F)}"]\\
h^*(F)\arrow[r,"h^*(S_g)"']&h^*(g^*(F)).
\end{tikzcd}
\]
这里 $C_{h,g}(F):(hg)^*(F)\longrightarrow(h^*g^*)(F)$ 是\Term{canonical}{典范}\Term{isomorphism}{同构}。

$X$ 上两个以 $G$ 为\Term{operator-group}{算子群}的\Term{sheaf}{层} $(F,S_g)$、$(F',S'_g)$ 之间的\Term{morphism}{态射}按显然的方式定义。由此得到\Term{category}{范畴} $\widetilde X_{\mathrm{ét},G}$；由 Giraud 判据（SGA 4 II 4.12）容易知道，它是一个\Term{topos}{拓扑斯}。

给定 $\widetilde X_{\mathrm{ét},G}$ 中的一个环 $A$，就可考虑 $A$-\Term{module}{模}的\Term{abelian-category}{阿贝尔范畴}，记为 $\operatorname{Mod}(A)$。

例如，给定一个环 $\Lambda$，利用\Term{canonical}{典范}的传递\Term{isomorphism}{同构}
\[
S_g:\Lambda_X\xrightarrow{\sim}g^*(\Lambda_X),
\]
可以在与 $\Lambda$ 对应的\Term{constant-sheaf}{常值层} $\Lambda_X$ 上赋予一个平凡的、以 $G$ 为\Term{operator-group}{算子群}的\Term{sheaf}{层}结构。

[T] 第 V 章所证明的一切结果都可推广到这个框架中，细节留给读者。

\SourcePage{386}
\SGAsection{2}{证明某个 $\Lambda[G]$-模复形是完美的}
先回顾（SGA 4 XVII）中的下述结果。
\begin{SGAresult}{\Term{lemma}{引理}}{2.1}
设 $X,Y$ 是两个\Term{scheme}{概形}，$f:X\longrightarrow Y$ 是\Term{proper-morphism}{固有态射}，$Y$ \Term{quasi-compact}{拟紧}。设 $\Lambda$ 是一个\Term{torsion-ring}{挠环}，$F^\bullet$ 是 $X$ 上具有\Term{finite}{有限}\Term{tor-dimension}{Tor 维数}（SGA 4 XVII）的 $\Lambda$-\Term{module}{模}\Term{sheaf}{层}\Term{complex}{复形}。则 $\mathrm Rf_*(F^\bullet)$ 具有\Term{finite}{有限}\Term{tor-dimension}{Tor 维数}。
\end{SGAresult}
\begin{SGAresult}{命题}{2.2}
设 $X$ 是\Term{separably-closed-field}{可分闭域} $k$ 上的\Term{proper-scheme}{固有概形}，$G$ 是\Term{action}{作用}于 $X$ 的\Term{finite-group}{有限群}，$V$ 是 $X$ 的一个在 $G$ 下\Term{stable}{稳定}的\Term{open-subset}{开集}，且 $G$ 在 $V$ 上的\Term{action}{作用}是\Term{admissible-action}{可容许的}（SGA I V 1.7）并且\Term{free-action}{自由}。设 $\Lambda=\mathbf Z/\ell^n\mathbf Z$，其中素数 $\ell$ 不同于 $k$ 的\Term{characteristic}{特征}；$\Lambda_V$ 是 $V$ 上与 $\Lambda$ 对应的\Term{constant-sheaf}{常值层}，带有其平凡的、以 $G$ 为\Term{operator-group}{算子群}的\Term{sheaf}{层}结构。记 $i:V\hookrightarrow X$ 为\Term{immersion}{浸入}，令 $\Lambda_{V,X}=i_!(\Lambda_V)$。则 $\Lambda[G]$-\Term{module}{模}\Term{complex}{复形} $\mathrm R\Gamma_X(\Lambda_{V,X})$ 是\Term{perfect-complex}{完美复形}（第 VIII 讲 7）。
\end{SGAresult}
\begin{proof}
设 $p:X\longrightarrow Y=X/G$ 为\Term{canonical-projection}{典范投影}。众所周知，在\Term{category}{范畴} $\mathrm D(\Lambda[G])$ 中有\Term{canonical}{典范}\Term{isomorphism}{同构}
\[
\mathrm R\Gamma_X(\Lambda_{V,X})\simeq
\mathrm R\Gamma_X(\mathrm Rp_*(\Lambda_{V,X})).
\]
另一方面，$Y$ 上的 $\Lambda[G]$-\Term{module}{模}\Term{complex}{复形} $\mathrm Rp_*(\Lambda_{V,X})$ 具有\Term{finite}{有限}\Term{tor-dimension}{Tor 维数}。事实上，$p:X\longrightarrow Y$ 是\Term{scheme}{概形}的\Term{finite-morphism}{有限态射}，所以 $q\ne0$ 时有 $\mathrm R^qp_*(\Lambda_{V,X})=0$（SGA 4 VIII 5.5）。因此
\SourcePage{387}
\[
\mathrm Rp_*(\Lambda_{V,X})=p_*(\Lambda_{V,X}).
\]
（SGA 4 VIII 5.5）还使我们能计算\Term{sheaf}{层} $p_*(\Lambda_{V,X})$ 的\Term{stalk}{茎}。考虑到 $G$ 在 $V$ 上\Term{free-action}{自由}\Term{action}{作用}，得到
\[
p_*(\Lambda_{V,X})_{\bar y}\simeq
\begin{cases}
\displaystyle\prod_{x\in p^{-1}(y)}\Lambda_{V,X,x}\simeq\Lambda[G],&y\in p(V),\\
0,&y\notin p(V),
\end{cases}
\]
其中 $\bar y$ 是位于 $y$ 上的 $Y$ 的\Term{geometric-point}{几何点}。

所以 $\mathrm Rp_*(\Lambda_{V,X})=p_*(\Lambda_{V,X})$ 在 $\Lambda[G]$ 上\Term{flat}{平坦}，因为它逐\Term{stalk}{茎}\Term{flat}{平坦}（SGA 4 XVII），因而具有\Term{finite}{有限}\Term{tor-dimension}{Tor 维数}。对\Term{canonical}{典范}\Term{morphism}{态射} $Y\longrightarrow\operatorname{Spec}(k)$ 应用 2.1，得知
\[
\mathrm R\Gamma_X(\Lambda_{V,X})=
\mathrm R\Gamma_Y(\mathrm Rp_*(\Lambda_{V,X}))
\]
具有\Term{finite}{有限}\Term{tor-dimension}{Tor 维数}。

要断言 $\mathrm R\Gamma_X(\Lambda_{V,X})$ \Term{perfect-complex}{完美}，还须证明它\Term{pseudo-coherent}{伪凝聚}且\Term{cohomology}{上同调}\Term{bounded}{有界}（第 VIII 讲 8.1）。由于 $\Lambda[G]$ \Term{noetherian}{诺特}，要证明其\Term{pseudo-coherent}{伪凝聚}，只须证明各 $H^q(X,\Lambda_{V,X})$ 是\Term{finite-type-module}{有限型} $\Lambda[G]$-\Term{module}{模}（第 VIII 讲 7.2），\Term{equivalence}{等价}地说，是\Term{finite-type-module}{有限型} $\Lambda$-\Term{module}{模}；这由对 $\Lambda$-\Term{module}{模} $\Lambda_{V,X}$ 应用（SGA 4 XIV 1.1）得到。
\end{proof}

当 $X$ 是\Term{curve}{曲线}时，Weil 公式 6.1 将用已知的整体\Term{invariants}{不变量}及第 5 节将定义的\Term{local-invariant}{局部不变量}，给出类
$\operatorname{cl}_{K^\bullet(\Lambda[G])}(\mathrm R\Gamma_X(\Lambda_{V,X}))$ 的表达式。

\SourcePage{388}
\SGAsection{3}{有限群线性表示的回顾}
本节中，$\Lambda$ 表示一个\Term{commutative-ring}{交换环}，$G$ 表示一个群。

\textbf{3.1。} 考虑 $\operatorname{Mod}(\Lambda[G])$ 的\Term{full-subcategory}{满子范畴} $\mathscr C$，其\Term{object}{对象}是底层 $\Lambda$-\Term{module}{模}为\Term{finite-projective-module}{有限型投射模}的 $\Lambda[G]$-\Term{module}{模} $M$。对这样的 $M$ 和 $g\in G$，考虑 $\Lambda$-\Term{module}{模}的\Term{endomorphism}{自同态}
\[
g_M:M\longrightarrow M,\qquad g_M(x)=gx.
\]
由此得到函数
\[
\chi_M:G\longrightarrow\Lambda,\qquad
\chi_M(g)=\operatorname{Tr}_\Lambda^*(g),
\]
它是 $G$ 上的\Term{central-function}{中心函数}，称为 $\Lambda[G]$-\Term{module}{模} $M$ 的\Term{character}{特征标}。

将每个 $M\in\operatorname{Ob}(\mathscr C)$ 对应到其\Term{character}{特征标} $\chi_M$，就得到一个\Term{additive-function}{可加函数}（第 VIII 讲 1）。利用\Term{perfect-complex}{完美复形}\Term{endomorphism}{自同态}的迹的\Term{form}{形式}性质，甚至可在那些作为 $\Lambda$-\Term{module}{模}\Term{complex}{复形}是\Term{perfect-complex}{完美}的 $\Lambda[G]$-\Term{module}{模}\Term{complex}{复形}所成的\Term{category}{范畴}上得到一个\Term{additive-function}{可加函数} $\operatorname{Tr}$。因而得到\Term{homomorphism}{同态}
\[
R^\bullet_\Lambda(G)\xrightarrow{\operatorname{Tr}}\mathscr F(G,\Lambda),
\]
其中 $R^\bullet_\Lambda(G)$ 是 $G$ 在 $\Lambda$ 上的\Term{representation-ring}{表示环}（第 VIII 讲 10），$\mathscr F(G,\Lambda)$ 是\Term{central-function}{中心函数} $G\longrightarrow\Lambda$ 所成的\Term{module}{模}。当 $G$ \Term{finite}{有限}时，\Term{canonical}{典范}\Term{homomorphism}{同态} $K^\bullet(\Lambda[G])\longrightarrow R^\bullet_\Lambda(G)$ 与 $\operatorname{Tr}$ 的复合记为 $\operatorname{tr}$。

回顾下述结果（第 IX 讲\Term{theorem}{定理} 1、推论 3）。
\SourcePage{389}
\begin{SGAresult}{命题}{3.2}
设 $\mathbf Z_\ell$ 为 $\ell$ 进\Term{ring-of-integers}{整数环}，$G$ 为\Term{finite-group}{有限群}。则\Term{homomorphism}{同态}
\[
\operatorname{tr}:K^\bullet(\mathbf Z_\ell[G])\longrightarrow\mathscr F(G,\mathbf Z_\ell)
\]
是\Term{monomorphism}{单态射}。两个在 $\mathscr F(G,\mathbf Z_\ell)$ 中有相同迹函数的\Term{projective-module}{投射} $\mathbf Z_\ell[G]$-\Term{module}{模}彼此\Term{isomorphism}{同构}。
\end{SGAresult}

\textbf{3.3。} 设 $\alpha:G\longrightarrow G'$ 是群\Term{homomorphism}{同态}。\Term{restriction}{限制}\Term{operator}{算子}给出一个\Term{exact-functor}{正合函子}
\[
J:\operatorname{Mod}(\Lambda[G'])\longrightarrow\operatorname{Mod}(\Lambda[G]).
\]
它有\Term{left-adjoint}{左伴随}
\[
I:\operatorname{Mod}(\Lambda[G])\longrightarrow\operatorname{Mod}(\Lambda[G']),
\qquad I(F)=\Lambda[G']\otimes_{\Lambda[G]}F.
\]
若 $F$ \Term{finite-projective-module}{有限型投射}（相应地，\Term{finite-type-module}{有限型}），则 $I(F)$ 也如此。

若 $\varphi\in\mathscr F(G',\Lambda)$ 是 $G'$ 上的\Term{central-function}{中心函数}，则 $\varphi\circ\alpha$ 是 $G$ 上的\Term{central-function}{中心函数}，故有映射
\[
\alpha_*:\mathscr F(G',\Lambda)\longrightarrow\mathscr F(G,\Lambda).
\]
若 $\varphi$ 是 $\mathscr C(G')$ 中 $\Lambda[G']$-\Term{module}{模} $M$ 的\Term{character}{特征标}，则 $\alpha_*(\varphi)$ 是 $\Lambda[G]$-\Term{module}{模} $J(M)$ 的\Term{character}{特征标}。

现在假设 $\alpha:G\longrightarrow G'$ 为\Term{injectivity}{单射}，且 $G$ 在 $G'$ 中的指数\Term{finite}{有限}。若 $F$ 是 $\mathscr C(G)$ 中的 $\Lambda[G]$-\Term{module}{模}，其\Term{character}{特征标}为 $\psi:G\longrightarrow\Lambda$，则 $I(F)$ 属于 $\mathscr C(G')$，而 $I(F)$ 的\Term{character}{特征标} $\psi^*:G'\longrightarrow\Lambda$ 是由 $\psi$ \Term{induced-character}{诱导的特征标}，也就是由下式确定的 $G'$ 上的函数：
\SourcePage{390}
\[
\psi^*(g)=\sum_{t\in G'/G}\psi(t^{-1}gt),\tag{3.3.1}
\]
约定当 $t^{-1}gt\notin G$ 时 $\psi(t^{-1}gt)=0$。

\textbf{3.4。} 设 $X^\bullet,Y^\bullet$ 是两个 $A$-\Term{module}{模}\Term{complex}{复形}。与\Term{module}{模}的情形相同，若 $X^\bullet,Y^\bullet$ 中有一个另带 $A$-$B$-\Term{bimodule}{双模}结构，就能在 $\operatorname{Hom}_A^\bullet(X^\bullet,Y^\bullet)$ 上赋予一个 $B$-\Term{module}{模}\Term{complex}{复形}结构，因而也在 $\mathrm R\operatorname{Hom}_A(X^\bullet,Y^\bullet)$ 上赋予相应结构。

例如，假定 $X^\bullet$ 是 $\mathrm D^-(A)$ 的\Term{object}{对象}，$Y^\bullet$ 是 $A$-$B$-\Term{bimodule}{双模}\Term{complex}{复形}；确切地说，处于\Term{signature}{符号} ${}_{A-B}Y^\bullet$ 所表示的情形。特别地，$Y^\bullet$ 是 $\mathrm D(A)$ 的\Term{object}{对象}。取 $X^\bullet$ 的一个\Term{projective-resolution}{投射消解} $P^\bullet\longrightarrow X^\bullet$。那么 $\operatorname{Hom}_A^\bullet(P^\bullet,Y^\bullet)$ 显然是一个左 $B$-\Term{module}{模}\Term{complex}{复形}，即 $\mathrm D(B)$ 中的\Term{object}{对象}。若另用一个\Term{projective-resolution}{投射消解} $P_1^\bullet\longrightarrow X^\bullet$，所得到的 $B$-\Term{module}{模}\Term{complex}{复形} $\operatorname{Hom}_A^\bullet(P_1^\bullet,Y^\bullet)$ 在 $\mathrm D(B)$ 中与 $\operatorname{Hom}_A^\bullet(P^\bullet,Y^\bullet)$ \Term{canonical}{典范}\Term{isomorphism}{同构}。所以可以把 $\mathrm R\operatorname{Hom}_A(X^\bullet,Y^\bullet)$ 定义为 $\mathrm D(B)$ 中在唯一\Term{isomorphism}{同构}意义下确定的\Term{object}{对象}，它对 $\mathrm D^-(A)$ 中的 $X^\bullet$ 和 $\mathrm D(A\otimes_{\mathbf Z}B)$ 中的 $Y^\bullet$ 具有\Term{functoriality}{函子性}。

类似的考虑也适用于 $X^\bullet\otimes_A^{\mathrm L}Y^\bullet$。例如，若 $Y^\bullet$ 如上是 $A$-$B$-\Term{bimodule}{双模}\Term{complex}{复形}，$X^\bullet$ 是右\Term{bounded}{有界}的右 $A$-\Term{module}{模}\Term{complex}{复形}，则可将 $X^\bullet\otimes_A^{\mathrm L}Y^\bullet$ 定义为 $\mathrm D(B)$ 中在唯一\Term{isomorphism}{同构}意义下确定的\Term{object}{对象}。

\textbf{3.5。} 设 $\Lambda$ 为\Term{commutative-ring}{交换环}，$A,B$ 是两个 $\Lambda$-\Term{algebra}{代数}，
\SourcePage{391}
$C=A\otimes_\Lambda B$。因此，$C$-\Term{module}{模}就是 $A$-$B$-\Term{bimodule}{双模}。

若 $\Delta$ 是 $\mathrm D(A)$ 的一个\Term{full-triangulated-subcategory}{满三角子范畴}，则用 $\mathrm D(C)_\Delta$ 表示 $\mathrm D(C)$ 的如下\Term{full-triangulated-subcategory}{满三角子范畴}：它由那些视为 $\mathrm D(A)$ 的\Term{object}{对象}时属于 $\Delta$ 的\Term{complex}{复形}生成。

\begin{SGAresult}{命题}{3.6}
沿用 3.4、3.5 的记号，设 $\Delta$ 是 $\mathrm D(A)$ 的一个对\Term{direct-summand}{直和因子}\Term{stable}{稳定}的\Term{full-triangulated-subcategory}{满三角子范畴}。若 $P^\bullet$ 是 $\mathrm D(B)_{\mathrm{parf}}$ 的\Term{object}{对象}（第 VIII 讲 7），$M^\bullet$ 是 $\mathrm D(C)_\Delta$ 的\Term{object}{对象}，则 4.1 中定义的 $A$-\Term{module}{模}\Term{complex}{复形} $\mathrm R\operatorname{Hom}_B(P^\bullet,M^\bullet)\in\mathrm D(A)$ 属于 $\Delta$。特别地，得到\Term{triangulated-category}{三角范畴}的\Term{pairing}{配对}
\[
\mathrm R\operatorname{Hom}:\mathrm D(B)_{\mathrm{parf}}\times\mathrm D(C)_\Delta\longrightarrow\Delta,\tag{3.6.1}
\]
也就是对每个变元都\Term{exact-functor}{正合}的\Term{bifunctor}{双函子}，因而（第 VIII 讲 3）得到\Term{bilinear-map}{双线性映射}
\[
K^\bullet(B)\times K^\bullet(\mathrm D(C)_\Delta)
\xrightarrow{\operatorname{hom}}K(\Delta),\tag{3.6.2}
\]
满足
\[
K(a)\bigl(\operatorname{cl}_{K^\bullet(B)}(P),
\operatorname{cl}_{K(\mathrm D(C)_\Delta)}(M^\bullet)\bigr)
=\operatorname{cl}_{K(\Delta)}\bigl(\mathrm R\operatorname{Hom}_B(P^\bullet,M^\bullet)\bigr).
\tag{3.6.3}
\]
若将 $B$ 换为其\Term{opposite-ring}{反环} $B^\circ$，并将 $\mathrm R\operatorname{Hom}_B(P,M)$ 换为 $P^\bullet\otimes_B^{\mathrm L}M^\bullet$，也得到类似陈述。
\end{SGAresult}
\begin{proof}
先在 $P^\bullet$ 只有一个非零\Term{set-component}{分量}的特殊情形证明 $\mathrm R\operatorname{Hom}_B(P,M)$ 是 $\Delta$ 中的\Term{object}{对象}，再对 $P^\bullet$ 的非零\Term{set-component}{分量}数作归纳。
\end{proof}
\SourcePage{392}
\textbf{3.7。} 设 $P$ 为 $\Lambda[G]$-\Term{module}{模}，特别地也是 $\Lambda$-\Term{module}{模}。令 $\check P=\operatorname{Hom}_\Lambda(P,\Lambda)$；对每个 $u\in\operatorname{Hom}_\Lambda(P,\Lambda)$ 和 $g\in G$，规定
$(gu)(x)=u(g^{-1}x)$，就可在 $\check P$ 上赋予左 $\Lambda[G]$-\Term{module}{模}结构，即
\[
g_{\check P}={}^t(g^{-1})_P.\tag{3.7.1}
\]
这样得到的 $\Lambda[G]$-\Term{module}{模} $\check P$ 称为 $P$ 的\Term{dual-module}{对偶模}。

因此可以把 $\Lambda[G]$-\Term{module}{模}\Term{complex}{复形} $P^\bullet$ 的\Term{dual-complex}{对偶复形} $\check P^\bullet=\operatorname{Hom}_\Lambda^\bullet(P^\bullet,\Lambda)$ 定义为一个 $\Lambda[G]$-\Term{module}{模}\Term{complex}{复形}。

若 $\chi_P$ 为 $P$ 的\Term{character}{特征标}，则其\Term{dual-module}{对偶}的\Term{character}{特征标} $\chi_{\check P}$ 满足
\[
\chi_{\check P}(g)=\chi_P(g^{-1}),\tag{3.7.2}
\]
由（4.3.1）并利用公式
\[
\operatorname{Tr}_\Lambda^*(v)=\operatorname{Tr}_\Lambda^*({}^tv)
\]
立即得到这一点。

若 $P$ 是\Term{finite-projective-module}{有限型投射} $\Lambda[G]$-\Term{module}{模}，则 $\check P$ 显然也是\Term{finite-projective-module}{有限型投射} $\Lambda[G]$-\Term{module}{模}。此外，此时\Term{canonical}{典范}的 $\Lambda$-\Term{module}{模}\Term{homomorphism}{同态}
\[
P\longrightarrow\check{\check P}
\]
实际上是 $\Lambda[G]$-\Term{module}{模}的\Term{isomorphism}{同构}。这样就在 $K^\bullet(\Lambda[G])$ 上定义了一个\Term{involution}{对合}
\[
x\longmapsto\check x,
\]
它把 $P$ 的类变为 $\check P$ 的类。显然，若 $\mu:\Lambda\longrightarrow\Gamma$ 是\Term{commutative-ring}{交换环}\Term{homomorphism}{同态}，$u:\Lambda[G]\longrightarrow\Gamma[G]$ 为相应的\Term{homomorphism}{同态}，则
\[
u^*:K^\bullet(\Lambda[G])\longrightarrow K^\bullet(\Gamma[G])
\]
与 $K^\bullet(\Lambda[G])$ 和 $K^\bullet(\Gamma[G])$ 上的\Term{involution}{对合} $x\longmapsto\check x$ 交换。

\SourcePage{393}
设 $A$ 为 $\Lambda$-\Term{algebra}{代数}，并令
\[
A[G]=\Lambda[G]\otimes_\Lambda A.
\]
给定一个 $A[G]$-\Term{module}{模}\Term{complex}{复形} $X^\bullet$，用 $X^{\bullet G}$ 表示逐项取 $G$-\Term{invariant}{不变量}所得的 $A$-\Term{module}{模}\Term{complex}{复形}。

\begin{SGAresult}{命题}{3.8}
设 $G$ 是\Term{finite-group}{有限群}，$P^\bullet$ 是各\Term{set-component}{分量}均为\Term{finite-projective-module}{有限型投射模}的\Term{bounded}{有界} $\Lambda[G]$-\Term{module}{模}\Term{complex}{复形}，$M^\bullet$ 是 $A[G]$-\Term{module}{模}\Term{complex}{复形}。则有 $A$-\Term{module}{模}\Term{complex}{复形}的\Term{canonical}{典范}\Term{isomorphism}{同构}
\[
\operatorname{Hom}_{\Lambda[G]}^\bullet(P^\bullet,M^\bullet)
\simeq\operatorname{Hom}_\Lambda^\bullet(P^\bullet,M^\bullet)^G
\simeq(\check P^\bullet\otimes_\Lambda M^\bullet)^G
\simeq\check P^\bullet\otimes_{\Lambda[G]}M^\bullet.
\tag{3.8.1}
\]
\end{SGAresult}
证明使用\Term{module}{模}情形的熟知\Term{isomorphism}{同构}，留给读者。注意，第二项和第三项中 $G$ 在 $\operatorname{Hom}_\Lambda^\bullet(P^\bullet,M^\bullet)$ 和 $\check P^\bullet\otimes_\Lambda M^\bullet$ 上的\Term{action}{作用}，分别由它在 $P^\bullet,M^\bullet$ 以及 $\check P^\bullet,M^\bullet$ 上的\Term{action}{作用}通过结构搬移来定义。

\begin{SGAresult}{推论}{3.9}
若 $P^\bullet\in\operatorname{Ob}(\mathrm D(\Lambda[G])_{\mathrm{parf}})$，$M^\bullet\in\mathrm D^+(A[G])$，则在\Term{derived-category}{导出范畴} $\mathrm D(A)$ 中有下列\Term{isomorphism}{同构}：
\[
\mathrm R\operatorname{Hom}_{\Lambda[G]}(P^\bullet,M^\bullet)
=\mathrm R\Gamma^G(\check P^\bullet\otimes_\Lambda^{\mathrm L}M^\bullet)
\simeq\check P^\bullet\otimes_{\Lambda[G]}^{\mathrm L}M^\bullet.
\tag{3.9.1}
\]
\end{SGAresult}
\begin{proof}
由假设 $P^\bullet\in\mathrm D(\Lambda[G])_{\mathrm{parf}}$，可设 $P^\bullet$ \Term{bounded}{有界}且各\Term{set-component}{分量}都是\Term{finite-projective-module}{有限型投射模}。
\SourcePage{394}
所以
\[
\mathrm R\operatorname{Hom}_{\Lambda[G]}(P^\bullet,M^\bullet)
=\operatorname{Hom}_{\Lambda[G]}^\bullet(P^\bullet,M^\bullet),
\]
可以使用\Term{isomorphism}{同构}（4.3.4）。但 $\check P^\bullet$ 与 $P^\bullet$ 一样，也是各\Term{set-component}{分量}为\Term{finite-projective-module}{有限型投射模}的\Term{bounded}{有界}\Term{complex}{复形}，所以
\[
\check P^\bullet\otimes_{\Lambda[G]}M^\bullet\simeq
\check P^\bullet\otimes_{\Lambda[G]}^{\mathrm L}M^\bullet.
\]
同样，$\check P^\bullet\otimes_\Lambda M^\bullet=
\check P^\bullet\otimes_\Lambda^{\mathrm L}M^\bullet$。此外，\Term{complex}{复形} $\check P^\bullet\otimes_\Lambda M^\bullet$ 的各\Term{set-component}{分量}都是\Term{cohomologically-trivial}{上同调平凡的}（[3] IX 3），故有
\[
\mathrm R\Gamma^G(\check P^\bullet\otimes_\Lambda^{\mathrm L}M^\bullet)
\simeq(\check P^\bullet\otimes_\Lambda M^\bullet)^G.
\]
\end{proof}

\SGAsection{4}{Swan 表示}
设 $E$ 是\Term{algebraically-closed-field}{代数闭域} $k$ 上的一条\Term{smooth}{光滑}\Term{connected}{连通}\Term{algebraic-curve}{代数曲线}。若 $x\in E$ 是\Term{closed-point}{闭点}，则其\Term{local-ring}{局部环} $\mathscr O_{E,x}$（或简记 $\mathscr O_x$）是\Term{discrete-valuation-ring}{离散赋值环}，其\Term{fraction-field}{分式域}与 $E$ 的\Term{rational-function-field}{有理函数域}\Term{isomorphism}{同构}，其\Term{residue-field}{剩余域}是 $k$。相应的\Term{valuation}{赋值}记作 $v_x$。

\begin{SGAresult}{\Term{lemma}{引理}}{4.1}
设 $E$ 是\Term{algebraically-closed-field}{代数闭域} $k$ 上的\Term{smooth}{光滑}\Term{algebraic-curve}{代数曲线}，$g:E\longrightarrow E$ 是不同于\Term{identity-map}{恒等映射}的\Term{endomorphism}{自同态}。假定 $x\in E$ 是 $g$ 的一个孤立\Term{fixed-point}{不动点}，$\pi$ 是 $\mathscr O_{E,x}$ 的\Term{uniformizer}{素元}。若 $v_x$ 是 $\mathscr O_{E,x}$ 的相应\Term{valuation}{赋值}，则映射 $g$ 的\Term{fixed-point}{不动点} $x$ 的\Term{multiplicity}{重数} $\nu_x(g)$（第 IV 讲\SGAsourcegap{原 PDF 第 394 页在“IV”之后留有空白，缺少所引节号。}）等于 $v_x(g(\pi)-\pi)$。
\end{SGAresult}
\SourcePage{395}
\begin{proof}
在乘积 $E\times E$ 中，分别用 $\Delta$ 和 $\Gamma_g$ 表示\Term{diagonal}{对角线}及\Term{endomorphism}{自同态} $g$ 的\Term{graph}{图像}。设相应的\Term{immersion}{浸入}为
\[
\gamma:E\longrightarrow E\times E,\qquad
\gamma':E\longrightarrow E\times E.
\]
在 $x$ 及 $(x,x)$ 处的\Term{complete}{完备}\Term{local-ring}{局部环}上，它们由\Term{homomorphism}{同态}
\[
\gamma^*,\gamma'^*:A\widehat\otimes A\longrightarrow A,
\qquad A=\widehat{\mathscr O}_x,
\]
给出，其中
\[
\gamma^*(a\otimes b)=ab,\qquad
\gamma'^*(a\otimes b)=ag^*(b).
\]
要计算的是 $E\times E$ 中交 $\Delta\cdot\Gamma_g$ 在点 $\delta(x)$ 的\Term{multiplicity}{重数} $\nu_x(g)$。而\Term{divisor}{除子} $\Delta$ 对应于 $A\widehat\otimes A$ 中由 $1\otimes\pi-\pi\otimes1$ 生成的\Term{ideal}{理想}，所以 $\delta(x)$ 在交 $\Delta\cdot\Gamma_g$ 中的\Term{multiplicity}{重数}为
\[
\nu_x(g)=v_x\bigl(\gamma'^*(1\otimes\pi-\pi\otimes1)\bigr)
=v_x(g(\pi)-\pi).
\]
\end{proof}

\textbf{4.2。} 设 $C$ 是\Term{algebraically-closed-field}{代数闭域} $k$ 上的\Term{smooth}{光滑}、\Term{proper-scheme}{固有}、\Term{connected}{连通}\Term{algebraic-curve}{代数曲线}。设 $\eta$ 是其\Term{generic-point}{泛点}，记 $K=k(\eta)$ 为其\Term{rational-function-field}{有理函数域}。设
\[
K'\supset K,\qquad[K':K]=n
\]
为\Term{degree}{次数} $n$ 的\Term{finite}{有限}\Term{galois-extension}{Galois 扩张}，其\Term{galois-group}{Galois 群}为 $G$。记 $C'$ 为 $C$ 在 $K'$ 中的\Term{normalization}{正规化}，令
\SourcePage{396}
\[
p:C'\longrightarrow C
\]
为\Term{canonical-projection}{典范投影}。\Term{galois-group}{Galois 群} $G$ 显然\Term{action}{作用}于 $C'$。$y'\in C'$ 在 $G$ 中的\Term{stabilizer}{稳定子}记为 $G_{y'}$。

若 $\pi$ 是 $\mathscr O_{C',y'}$ 的\Term{uniformizer}{素元}，$s\in G_{y'}$ 且 $s\ne e$，令
\[
v_{y'}(s(\pi)-\pi)=\nu_{y'}(s).\tag{4.2.1}
\]
这就是映射 $s:C'\longrightarrow C'$ 的\Term{fixed-point}{不动点} $y'$ 的\Term{multiplicity}{重数}（\Term{lemma}{引理} 5.1）。

与 $y'$ 对应的\Term{swan-representation}{Swan 表示}由其\Term{character}{特征标} $\sigma_{y'}$ 定义，后者满足
\[
\sigma_{y'}(s)=
\begin{cases}
1-\nu_{y'}(s),&s\ne e,\\
\displaystyle\sum_{s\ne e}(\nu_{y'}(s)-1),&s=e.
\end{cases}\tag{4.2.2}
\]
Artin--Serre--Swan 的结果（第 IX 讲\Term{theorem}{定理} 4）表明，若 $\ell$ 为不同于 $k$ 的\Term{characteristic}{特征}的素数，则存在一个\Term{finite-projective-module}{有限型投射} $\mathbf Z_\ell[G_{y'}]$-\Term{module}{模} $\operatorname{Sw}_{y'}$，它在\Term{isomorphism}{同构}意义下确定，其\Term{character}{特征标}为 $\sigma_{y'}$。

还值得考虑 $\mathbf Z_\ell[G_{y'}]$-\Term{module}{模}
\[
\operatorname{Sw}'_{y'}=\operatorname{Sw}_{y'}\oplus\mathbf Z_\ell[G_{y'}].\tag{4.2.3}
\]
$\operatorname{Sw}'_{y'}$ 的\Term{character}{特征标} $\sigma'_{y'}$ 满足
\SourcePage{397}
\[
\sigma'_{y'}(s)=
\begin{cases}
1-\nu_{y'}(s),&s\ne e,\\
\displaystyle1+\sum_{s\ne e}\nu_{y'}(s)=1+v_{y'}(\mathfrak D_{C'/C}),&s=e,
\end{cases}\tag{4.2.4}
\]
其中 $\mathfrak D_{C'/C}$ 是\Term{covering}{覆盖} $C'/C$ 的\Term{different}{不同理想}。最后一个等式就是 [3，第 VI 章第 2 节，命题 4] 中取 $H=\{1\}$ 的情形。

最后，函数
\[
a_{y'}=\sigma_{y'}+u_{G_{y'}}=\sigma'_{y'}-1_{G_{y'}}\tag{4.2.5}
\]
也是某个 $\mathbf Q_\ell[G_{y'}]$-\Term{module}{模}的\Term{character}{特征标}，这个\Term{module}{模}称为 $G_{y'}$ 的\Term{artin-representation}{Artin 表示}，记为 $\operatorname{Art}_{y'}$；这里 $u_{G_{y'}}$ 表示 $\mathbf Z_\ell[G]$ 的\Term{augmentation-ideal}{增广理想}的\Term{character}{特征标}。因此
\[
a_{y'}(s)=
\begin{cases}
-\nu_{y'}(s),&s\ne e,\\
\displaystyle\sum_{s\ne e}\nu_{y'}(s)=v_{y'}(\mathfrak D_{C/C'}),&s=e.
\end{cases}\tag{4.2.6}
\]
现在考虑包含 $G_{y'}\subset G$。令
\[
\operatorname{Sw}_y=\operatorname{Sw}_{y'}\otimes_{\mathbf Z_\ell[G_{y'}]}\mathbf Z_\ell[G],\qquad
\operatorname{Sw}'_y=\operatorname{Sw}'_{y'}\otimes_{\mathbf Z_\ell[G_{y'}]}\mathbf Z_\ell[G],\tag{4.2.7}
\]
\[
\operatorname{Art}_y=\operatorname{Art}_{y'}\otimes_{\mathbf Q_\ell[G_{y'}]}\mathbf Q_\ell[G].\tag{4.2.8}
\]
上述\Term{representation}{表示}的\Term{character}{特征标}分别记为 $\sigma_y,\sigma'_y,a_y$。

\SourcePage{398}
\textbf{4.3。} 注意，如此得到的\Term{module}{模}在\Term{isomorphism}{同构}意义下只依赖点 $y\in C$，不依赖于 $y$ 上方的 $C'$ 中的点 $y'$ 的选择。设 $y''$ 是另一个这样的点，则存在 $g\in G$ 使得 $y''=g\cdot y'$。于是，由 $C'$ 上的\Term{isomorphism}{同构} $g$ 和 $G$ 上的 $\operatorname{int}(g):h\longmapsto ghg^{-1}$，可把情形 $(C',y',G)$ 变为 $(C',y'',G)$。通过结构搬移可知，与 $y''$ 相应的\Term{module}{模} $\operatorname{Sw}_y$ 等，是由与 $y'$ 相应的\Term{module}{模}沿 $\operatorname{int}(g):G\longrightarrow G$ \Term{group-extension}{扩张}\Term{operator}{算子}得到的。我们知道，这种\Term{group-extension}{扩张}把一个\Term{module}{模}变为与之\Term{isomorphism}{同构}的\Term{module}{模}。

\begin{SGAresult}{命题}{4.4}
沿用 3.7 的记号，取 $\Lambda=\mathbf Z_\ell$，并设 $M^\bullet$ 为 $\mathrm D^+(\Lambda[G])$ 的\Term{object}{对象}。则有\Term{isomorphism}{同构}
\[
\operatorname{Sw}_y\simeq\check{\operatorname{Sw}}_y,\tag{*}
\]
\[
\operatorname{Hom}_{\mathbf Z_\ell[G]}^\bullet(\operatorname{Sw}_y,M^\bullet)
\simeq\operatorname{Sw}_y\otimes_{\mathbf Z_\ell[G]}M^\bullet.\tag{**}
\]
\end{SGAresult}
\begin{proof}
两个\Term{finite-projective-module}{有限型投射} $\mathbf Z_\ell[G]$-\Term{module}{模}若有相同\Term{character}{特征标}就\Term{isomorphism}{同构}（3.2）。若用 $\check\sigma_y$ 表示 $\check{\operatorname{Sw}}_y$ 的\Term{character}{特征标}，则由（3.7.2）有
\[
\check\sigma_y(s)=\sigma_y(s^{-1}).
\]
另一方面，
\[
v_{y'}(s^{-1}(\pi)-\pi)
=v_{y'}\bigl(s(s^{-1}(\pi)-\pi)\bigr)
=v_{y'}(\pi-s(\pi))=v_{y'}(s(\pi)-\pi),
\]
所以 $\check{\operatorname{Sw}}_{y'}=\operatorname{Sw}_{y'}$，由定义（4.2.7）即得公式（*）。再由 3.8 得到公式（**）。
\end{proof}

\SourcePage{399}
\textbf{4.5。} 下面要用到\Term{swan-representation}{Swan 表示}的\Term{character}{特征标}的一个显式公式。事实上有
\[
\sigma_{y'}^*(s)=\sigma_{y''}^*(s)=\sigma_y(s)=
\begin{cases}
\displaystyle\sum_{\substack{p(z')=y\\s(z')=z'}}(1-\nu_{z'}(s)),&s\ne e,\\
\displaystyle\sum_{p(z')=y}(1+v_{z'}(\mathfrak D_{C'/C}))-n,&s=e;
\end{cases}\tag{4.5.1}
\]
\[
\sigma_{y'}^{\prime *}(s)=\sigma_{y''}^{\prime *}(s)=\sigma'_y(s)=
\begin{cases}
\displaystyle\sum_{\substack{p(z')=y\\s(z')=z'}}(1-\nu_{z'}(s)),&s\ne e,\\
\displaystyle\sum_{p(z')=y}(1+v_{z'}(\mathfrak D_{C'/C})),&s=e;
\end{cases}\tag{4.5.2}
\]
\[
a_{y'}^*(s)=a_{y''}^*(s)=a_y(s)=
\begin{cases}
\displaystyle-\sum_{\substack{p(z')=y\\s(z')=z'}}\nu_{z'}(s),&s\ne e,\\
\displaystyle\sum_{p(z')=y}v_{z'}(\mathfrak D_{C'/C}),&s=e.
\end{cases}\tag{4.5.3}
\]
为\Term{derive}{导出}这些公式，使用（3.3.1），并注意：固定一点 $y'_0\in C'$ 后，对每个满足 $p(y')=p(y'_0)$ 的 $y'\in C'$，都可找到 $g'\in G$ 使 $g'(y')=y'_0$。这些 $g'$ 构成 $G$ 中 $G_{y'_0}$ 的右陪集的一组代表，且有 $G_{y'_0}=g'G_{y'}g'^{-1}$。在（3.3.1）的右边使用这组代表，就得到上述公式。

\SourcePage{400}
\textbf{4.6。} \Term{invariants}{不变量} $\operatorname{Sw}_y,\operatorname{Sw}'_y,\operatorname{Art}_y$ 可以在更一般的条件下定义。设 $C'$ 是 $k$ 上的一条\Term{smooth}{光滑}\Term{algebraic-curve}{代数曲线}，未必\Term{connected}{连通}。假定\Term{finite-group}{有限群} $G$ \Term{action}{作用}于 $C'$，且在 $C=C'/G$ 中存在一个\Term{dense-open}{稠密开集} $U$，使得 $C'|U$ 是 $U$ 上群为 $G$ 的\Term{principal-covering}{主覆盖}。对点 $y\in C$，按上述方法即可定义\Term{module}{模} $\operatorname{Sw}_y,\operatorname{Sw}'_y,\operatorname{Art}_y$。可以如下归约到 $C'$ \Term{connected}{连通}的情形。取满足 $p(y')=y$ 的 $y'\in C'$，设 $C'_1$ 是包含 $y'$ 的 $C'$ 的\Term{irreducible-component}{不可约分支}，$C_1$ 是包含 $y$ 的 $C$ 的\Term{irreducible-component}{不可约分支}。若 $G_{C'_1}$ 是分支 $C'_1$ 的\Term{stabilizer}{稳定子}，$\eta'_1$（相应地，$\eta_1$）为 $C'_1$（相应地，$C_1$）的\Term{generic-point}{泛点}，则\Term{group-extension}{扩张}
\[
K'_1=k(\eta'_1)\supset k(\eta_1)=K_1
\]
是\Term{galois-group}{Galois 群}为 $G_{C'_1}$ 的\Term{galois-extension}{Galois 扩张}，且 $C'_1$ 是 $C_1$ 在 $K'_1$ 中的\Term{normalization}{正规化}（SGA I I 10）。因此，可以对\Term{curve}{曲线} $C'_1,C_1$ 和群 $G_{C'_1}$ 应用上述讨论。由此得到\Term{finite-projective-module}{有限型投射} $\mathbf Z_\ell[G_{C'_1}]$-\Term{module}{模}
\[
\operatorname{Sw}_{G_{C'_1},y},\quad
\operatorname{Sw}'_{G_{C'_1},y},\quad
\operatorname{Art}_{G_{C'_1},y},
\]
而 $\operatorname{Sw}_y,\operatorname{Sw}'_y,\operatorname{Art}_y$ 就是相对于包含 $G_{C'_1}\subset G$ 的相应\Term{induced-module}{诱导模}。

\textbf{注记 4.6。} \Term{module}{模} $\operatorname{Sw}_{y'},\operatorname{Sw}'_{y'},\operatorname{Art}_{y'}$ 是纯\Term{local-scheme}{局部}的\Term{invariants}{不变量}：只要知道群 $G_{y'}$ 在\Term{local-ring}{局部环} $\mathscr O_{C',y'}$ 的\Term{completion}{完备化}上的\Term{action}{作用}，就确定了这些\Term{module}{模}。此外，$\operatorname{Sw}_{y'}=0$ 当且仅当 $p:C'\longrightarrow C$ 在 $y'$ 处\Term{tamely-ramified}{温和分歧}，也\Term{equivalence}{等价}于 $G_{y'}$ 的阶与 $k$ 的\Term{characteristic}{特征}互素 [3 VI]。
\SourcePage{401}
\SGAsection{5}{Weil 公式（参见 [3 VI §4]）}
\textbf{5.0。} 本节考虑的所有\Term{scheme}{概形}都定义在一个\Term{algebraically-closed-field}{代数闭}\Term{base-field}{基域} $k$ 上。

使用如下记号：$\Lambda_\nu=\mathbf Z/\ell^\nu\mathbf Z$，其中 $\ell$ 是给定的、不同于 $k$ 的\Term{characteristic}{特征}的素数。\Term{canonical}{典范}\Term{homomorphism}{同态} $\mathbf Z_\ell\longrightarrow\Lambda_\nu$ 诱导\Term{homomorphism}{同态}
\[
\lambda_\nu:\mathbf Z_\ell[G]\longrightarrow\Lambda_\nu[G],
\]
从而 [第 VIII 讲 9] 诱导\Term{reduced}{约化}\Term{homomorphism}{同态}
\[
\lambda_\nu^*:K^\bullet(\mathbf Z_\ell[G])\longrightarrow K^\bullet(\Lambda_\nu[G]).
\]
与\Term{module}{模} $\operatorname{Sw}_y,\operatorname{Sw}'_y$ 对应的元素在 $\lambda_\nu^*$ 下的像，分别记为 $\operatorname{sw}_y^{\Lambda_\nu},\operatorname{sw}_y^{\prime\Lambda_\nu}$。

\Term{reduced}{约化}\Term{homomorphism}{同态} $\lambda_\nu^*$ 实际上是\Term{isomorphism}{同构} [第 IX 讲 1.5]。它的逆记为 $\mu_\nu$。

\Term{scheme}{概形} $X$ 上与环 $\Lambda_\nu$ 对应的\Term{constant-sheaf}{常值层}记为 $\Lambda_{\nu,X}$。设 $V$ 为 $X$ 的\Term{open-subset}{开集}，$j:V\hookrightarrow X$ 为\Term{canonical}{典范}\Term{immersion}{浸入}。令
\[
j_!(\Lambda_{\nu,V})=\Lambda_{\nu,V,X}.
\]

\SourcePage{402}
\begin{SGAresult}{命题}{5.1}
设 $C'$ 是\Term{algebraically-closed-field}{代数闭域} $k$ 上的\Term{smooth}{光滑}\Term{proper-scheme}{固有}\Term{algebraic-curve}{代数曲线}，$G$ 为\Term{action}{作用}于 $C'$ 的\Term{finite-group}{有限群}，$C=C'/G$，$p:C'\longrightarrow C$ 为\Term{canonical-projection}{典范投影}。设 $U$ 为 $C$ 的\Term{dense-open}{稠密开集}，且 $U'=p^{-1}(U)$ 是 $U$ 上群为 $G$ 的\Term{principal-covering}{主覆盖}；令 $Y=C-U$，$Y'=p^{-1}(X)=C'-U'$。记 $\operatorname{EP}(C)$ 为 $C$ 的\Term{euler-poincare-characteristic}{Euler--Poincaré 特征数}（若 $C$ \Term{connected}{连通}且\Term{genus}{亏格}为 $g$，它等于 $2-2g$）。则有 Weil 公式
\[
\operatorname{cl}_{K^\bullet(\Lambda_\nu[G])}
\bigl(\mathrm R\Gamma_{C'}(\Lambda_{\nu,U',C'})\bigr)
=\operatorname{EP}(C)\operatorname{cl}_{K^\bullet(\Lambda_\nu[G])}(\Lambda_\nu[G])
-\sum_{y\in Y}\operatorname{sw}_y^{\prime\Lambda_\nu},\tag{5.2}
\]
其中（6.2）的左边由 2.1 得到定义。
\end{SGAresult}
\begin{proof}
分别用 $S_1,S_2$ 表示待证公式的左边和右边。\Term{reduced}{约化}\Term{homomorphism}{同态} $\lambda_\nu^*$ 实际上是\Term{isomorphism}{同构}，故只须证明 $\mu_\nu(S_1)=\mu_\nu(S_2)$。

另一方面（3.3.1），迹\Term{homomorphism}{同态}
\[
K^\bullet(\mathbf Z_\ell[G])\xrightarrow{\operatorname{tr}}\mathscr F(G,\mathbf Z_\ell)
\]
是\Term{monomorphism}{单态射}。所以只须证明
\[
\operatorname{tr}(\mu_\nu(S_1))=\operatorname{tr}(\mu_\nu(S_2)).
\]
\Term{character}{特征标} $\sigma_2=\operatorname{tr}(\mu_\nu(S_2))$ 显然由公式
\[
\sigma_2=\operatorname{EP}(C)r_G-\sum_{y\in Y}\sigma'_y\tag{5.2.1}
\]
\SourcePage{403}
给出，其中 $\sigma'_y$ 是 $\operatorname{Sw}'_y$ 的\Term{character}{特征标}（（4.2.7）和（4.2.3）），$r_G$ 是\Term{regular-representation}{正则表示}。

现在考虑迹的\Term{reduced}{约化}\Term{homomorphism}{同态}
\[
\mathscr F(G,\mathbf Z_\ell)\xrightarrow{\lambda_n^\natural}
\mathscr F(G,\mathbf Z/\ell^n\mathbf Z).
\]
显然有
\[
\mathscr F(G,\mathbf Z_\ell)=\varprojlim_n(G,\mathbf Z/\ell^n\mathbf Z).
\]
因此，只须证明对每个 $n$ 有
\[
\lambda_n^\natural\bigl(\operatorname{tr}(\mu_\nu(S_1))\bigr)
=\lambda_n^\natural(\sigma_2).
\]
左边等于 $\operatorname{tr}(\lambda_n^*(\mu_\nu(S_1)))$，而
\[
\lambda_n^*(\mu_\nu(S_1))=
\operatorname{cl}_{K^\bullet(\Lambda_n[G])}
\bigl(\mathrm R\Gamma_{C'}(\Lambda_{\nu,U',C'})\bigr),
\]
这是由于 \Term{kunneth-formula}{Künneth 公式}（SGA 4 XVIII）
\[
\mathrm R\Gamma_{C'}(\Lambda_{n,U',C'})\otimes_{\Lambda_n}^{\mathrm L}\Lambda_\nu
\simeq\mathrm R\Gamma_{C'}(\Lambda_{\nu,U',C'}),\qquad n\geq\nu.
\]
于是归结为证明
\[
\operatorname{tr}\left(\operatorname{cl}_{K^\bullet(\Lambda_n[G])}
\bigl(\mathrm R\Gamma_{C'}(\Lambda_{n,U',C'})\bigr)\right)
=\lambda_n^\natural(\sigma_2).
\]
用 $\tau$ 表示这个等式的左边，并取 $g\in G$。$g$ 显然诱导下述\Term{distinguished-triangle}{区分三角}的一个\Term{endomorphism}{自同态}：
\SourcePage{404}
\[
\begin{tikzcd}[column sep=large]
&\mathrm R\Gamma_{Y'}(\Lambda_{n,Y'})\arrow[dl]\arrow[dr]&\\
\mathrm R\Gamma_{C'}(\Lambda_{n,U',C'})\arrow[rr]&&\mathrm R\Gamma_{C'}(\Lambda_{n,C'}).
\end{tikzcd}
\]
因此
\[
\begin{aligned}
\tau(g)&=\operatorname{Tr}^*g_{\mathrm R\Gamma_{C'}(\Lambda_{n,U',C'})}\\
&=\operatorname{Tr}^*g_{\mathrm R\Gamma_{C'}(\Lambda_{n,C'})}
-\operatorname{Tr}^*g_{\mathrm R\Gamma_{Y'}(\Lambda_{n,Y'})}.
\end{aligned}
\]
为简便，写 $\operatorname{Tr}^*$ 代替 $\operatorname{Tr}_{\Lambda_n}^*$。对带有 $g$ \Term{action}{作用}的\Term{discrete-space}{离散空间} $Y'$ 应用一个平凡的“Lefschetz 公式”，得到
\[
\operatorname{Tr}^*g_{\mathrm R\Gamma_{Y'}(\Lambda_{n,Y'})}
=\text{$g$ 在 $Y'$ 上的限制 $g_{Y'}$ 的不动点数}
=\sum_{\substack{y'\in Y'\\gy'=y'}}1.
\]
另一方面，对\Term{action}{作用}于 $C'$ 的 $g$，Lefschetz 公式（第 III 讲 4.7）在 $g\ne e$ 时给出
\[
\operatorname{Tr}^*g_{\mathrm R\Gamma_{C'}(\Lambda_{n,Y'})}
=\sum_{y'\in Y'}\nu_{y'}(g)
=\text{映射 $g_{C'}:C'\longrightarrow C'$ 的不动点重数之和},
\]
而当 $g=e$ 时，有
\[
\operatorname{Tr}^*g_{\mathrm R\Gamma_{C'}(\Lambda_{n,C'})}
=\operatorname{EP}(C')\cdot1_\Lambda;
\]
这是由通过（SGA 4 IX 4.7）以及
\SourcePage{405}
\Term{kunneth-formula}{Künneth 公式}（SGA 4 XVII）对 $H^i(C',\Lambda_{n,C'})$ 的显式计算得到的。

由此得到 $\tau$ 的公式
\[
\tau(g)=
\begin{cases}
\displaystyle\sum_{\substack{y'\in Y'\\g(y')=y'}}(\nu_{y'}(g)-1),&g\ne e,\\
\operatorname{EP}(C')-\operatorname{card}(Y'),&g=e.
\end{cases}\tag{5.2.2}
\]
与（6.2.1）比较，容易得到 $\tau=\lambda_n^\natural(\sigma_2)$。事实上，若 $g\ne e$，由（4.5.2）立即得到 $\tau(g)=\lambda_n^\natural(\sigma_2)(g)$。若 $g=e$，要验证的公式为
\[
\operatorname{EP}(C')-\operatorname{card}(Y')
=N\operatorname{EP}(C)-\sum_{y\in Y}\sigma'_y(e),
\qquad N=\operatorname{card}(G),
\]
即（4.5.2）
\[
\operatorname{EP}(C')-N\operatorname{EP}(C)
=-\deg(\mathfrak D_{C'/C}),
\]
这正是经典的 Hurwitz 公式。
\end{proof}

\SGAsection{6}{局部项 $\varepsilon_x^\Delta(F)$ 的定义}
\textbf{6.1。} 设 $C$ 是\Term{algebraically-closed-field}{代数闭域} $k$ 上的\Term{connected}{连通}\Term{smooth}{光滑}\Term{algebraic-curve}{代数曲线}，$A$ 是 $\Lambda$-\Term{algebra}{代数}（$\Lambda=\mathbf Z/\ell^\nu\mathbf Z$，$\ell$ 不同于 $k$ 的\Term{characteristic}{特征}），$\Delta\subset\mathrm D^+(A)$ 是对\Term{direct-summand}{直和因子}\Term{stable}{稳定}的\Term{triangulated-subcategory}{三角子范畴}。设 $F$ 是 $C$ 上的 $A$-\Term{module}{模}\Term{sheaf}{层}\Term{complex}{复形}，满足以下条件：

i) 对每个\Term{geometric-point}{几何点} $\bar x\in C(k)$，有 $F_{\bar x}\in\operatorname{Ob}(\Delta)$。
\SourcePage{406}

ii) 存在 $C$ 的一个\Term{nonempty}{非空}\Term{open-subset}{开集}，使 $F$ 在其上“\Term{locally-constant}{局部常值}”；\Term{equivalence}{等价}地说，存在\Term{nonempty}{非空}\Term{open-subset}{开集} $U\subset C$ 及 $C$ 的一个\Term{connected}{连通}\Term{galois-covering}{Galois 覆盖} $p:C'\longrightarrow C$，其群为 $G$，使得若 $F'=p^*(F)$ 是 $F$ 沿 $p$ 的\Term{inverse-image}{逆像}，则 $F'|U'$ \Term{constant-sheaf}{常值}，其中 $U'=p^{-1}(U)$。

于是\Term{sheaf}{层} $F|U$ 由一个带 $G$ \Term{action}{作用}的 $A$-\Term{module}{模}\Term{complex}{复形} $M$ 定义。若分别用 $\eta,\eta'$ 表示 $C,C'$ 的\Term{generic-point}{泛点}，用 $\bar\eta,\bar\eta'$ 表示位于 $\eta,\eta'$ 上的\Term{geometric-point}{几何点}，则有
\[
p(\bar\eta')=\bar\eta,\qquad
F_{\bar\eta}\simeq F'_{\bar\eta'}=M.\tag{6.1.1}
\]
作为 $A$-\Term{module}{模}\Term{complex}{复形}，$M$ 与 $F_{\bar\eta}$ \Term{isomorphism}{同构}，由条件 i)，它是 $\Delta$ 的\Term{object}{对象}。所以（3.6），对每个 $x\in C(k)$，$\operatorname{Hom}_{\Lambda[G]}^\bullet(\operatorname{Sw}_x^\Lambda,M)$ 是 $\Delta$ 的\Term{object}{对象}，因而给出 $K(\Delta)$ 的一个元素。

事实上，由 $\operatorname{Hom}_{\Lambda[G]}^\bullet(\operatorname{Sw}_x^\Lambda,M)$ 定义的 $\Delta$ 的\Term{object}{对象}，在\Term{isomorphism}{同构}意义下不依赖于使 $F|U$ \Term{trivialization}{平凡化}的\Term{galois-covering}{Galois 覆盖} $p:C'\longrightarrow C$。要证明这一点，只须比较与两个\Term{covering}{覆盖}
\[
C'\xrightarrow{p}C,\qquad C''\xrightarrow{p'}C'\xrightarrow{p}C
\]
相应的元素。此时，若 $G'$ 是\Term{covering}{覆盖} $C''\longrightarrow C$ 的\Term{galois-group}{Galois 群}，则 $G$ 为 $G'$ 的\Term{quotient-group}{商群}。设
\[
\varphi:\mathbf Z_\ell[G']\longrightarrow\mathbf Z_\ell[G]
\]
为相应的\Term{canonical}{典范}\Term{epimorphism}{满同态}。相应的\Term{swan-module}{Swan 模}满足关系
\SourcePage{407}
\[
\operatorname{Sw}_{G,x}\simeq
\operatorname{Sw}_{G',x}\otimes_{\mathbf Z_\ell[G']}\mathbf Z_\ell[G],
\]
因为\Term{character}{特征标} $a_{G,x}$ 是从 $G'$ 到 $G$ 的\Term{induced-character}{诱导特征标}（[3] VI 命题 3），立即得到此式。

另一方面，有 $A[G']$-\Term{module}{模}\Term{complex}{复形}的\Term{isomorphism}{同构}
\[
F_{\bar\eta''}\simeq[F_{\bar\eta'}](\varphi).
\]
由此得到 $A$-\Term{module}{模}\Term{complex}{复形}
\[
\operatorname{Hom}_{\Lambda[G]}^\bullet(\operatorname{Sw}_{G,x}^\Lambda,F_{\bar\eta'})
\quad\text{与}\quad
\operatorname{Hom}_{\Lambda[G']}^\bullet(\operatorname{Sw}_{G',x}^\Lambda,F_{\bar\eta''})
\]
之间的\Term{canonical}{典范}\Term{isomorphism}{同构}（\Term{adjunction}{伴随}公式）。

\textbf{6.2。} 使用以下记号：
\[
\alpha_x^\Delta(F)=\operatorname{cl}_{K(\Delta)}
\left(\operatorname{Hom}_{\Lambda[G]}^\bullet(\operatorname{Sw}_x^\Lambda,F_{\bar\eta})\right),\tag{6.2.1}
\]
\[
\varepsilon_x^\Delta(F)=\alpha_x^\Delta(F)+
\bigl(\operatorname{cl}_{K(\Delta)}(F_{\bar\eta})-
\operatorname{cl}_{K(\Delta)}(F_{\bar x})\bigr).\tag{6.2.2}
\]
当然，在（6.2.1）中，是通过\Term{isomorphism}{同构}（6.1.1）在 $F_{\bar\eta}$ 上赋予 $A[G]$-\Term{module}{模}结构。注意，若 $x\in C(k)$，且 $C'$ 在 $x$ 处\Term{tamely-ramified}{温和分歧}，则由（4.6）有 $\operatorname{Sw}_x=0$，故 $\alpha_x^\Delta(F)=0$；此时（与 Ogg--Chafarevitch 一致）有
\[
\varepsilon_x^\Delta(F)=\operatorname{cl}_{K(\Delta)}(F_{\bar\eta})-
\operatorname{cl}_{K(\Delta)}(F_{\bar x}).\tag{6.2.3}
\]
显然，若 $x\notin Y=C-U$，则 $\varepsilon_x^\Delta(F)=0$。因此只有\Term{finite}{有限}多个 $x$ 满足 $\varepsilon_x^\Delta(F)\ne0$。

\SourcePage{408}
下面给出满足上述条件的\Term{category}{范畴} $\Delta$ 和\Term{sheaf}{层}\Term{complex}{复形} $F$ 的例子。

a) 取 $\Delta=\mathrm D^b(A)_{\mathrm{coh}}$（第 VIII 讲 5），$F$ 为\Term{constructible}{可构造}（SGA 4 XVII）且\Term{cohomology}{上同调}\Term{bounded}{有界}的 $A$-\Term{module}{模}\Term{sheaf}{层}\Term{complex}{复形}。此时 $\varepsilon_x^\Delta(F)\in K_\bullet(A)$。

b) 取 $\Delta=\mathrm D(A)_{\mathrm{parf}}$（第 VIII 讲 6），$F$ 为\Term{constructible}{可构造}且具有\Term{finite}{有限}\Term{tor-dimension}{Tor 维数}的 $A$-\Term{module}{模}\Term{sheaf}{层}\Term{complex}{复形}。此时 $\varepsilon_x^\Delta(F)\in K^\bullet(A)$。

c) 设 $A,B$ 为两个 $\Lambda$-\Term{algebra}{代数}，$\varphi:A\longrightarrow B$ 为 $\Lambda$-\Term{algebra}{代数}\Term{homomorphism}{同态}。取 $\Delta=\mathrm D^b(B\ \mathrm{rel}\ A)_{\mathrm{coh}}$（第 VIII 讲 10），$F$ 为 $B$-\Term{module}{模}\Term{sheaf}{层}\Term{complex}{复形}，它作为 $A$-\Term{module}{模}\Term{complex}{复形}的\Term{sheaf}{层}是\Term{constructible}{可构造}的，且\Term{cohomology}{上同调}\Term{bounded}{有界}。得到 $\varepsilon_x^\Delta(F)\in K_\bullet(B\ \mathrm{rel}\ A)$。显然，也可以像上面那样作变动，例如得到元素 $\varepsilon_x^\Delta(F)\in K^\bullet(B\ \mathrm{rel}\ A)$。

\SGAsection{7}{Euler--Poincaré 公式}
\begin{SGAresult}{\Term{theorem}{定理}}{7.1}
设 $C$ 是\Term{algebraically-closed-field}{代数闭域} $k$ 上的\Term{connected}{连通}、\Term{smooth}{光滑}、\Term{proper-scheme}{固有}\Term{algebraic-curve}{代数曲线}，$\eta$ 为 $C$ 的\Term{generic-point}{泛点}，$\Delta\subset\mathrm D^+(A)$ 是对\Term{direct-summand}{直和因子}\Term{stable}{稳定}的\Term{triangulated-subcategory}{三角子范畴}，$F$ 是 $A$-\Term{module}{模}\Term{sheaf}{层}\Term{complex}{复形}，满足 7.1 的条件 i)、ii)。则 $\mathrm R\Gamma_C(F)\in\operatorname{Ob}(\Delta)$，并且有下列公式（\Term{euler-poincare-formula}{Euler--Poincaré 公式}）：
\SourcePage{409}
\[
\operatorname{cl}_{K(\Delta)}(\mathrm R\Gamma_C(F))
=\operatorname{EP}(C)\operatorname{cl}_{K(\Delta)}(F_{\bar\eta})
-\sum_{x\in C(k)}\varepsilon_x^\Delta(F),\tag{7.2}
\]
其中 $\varepsilon_x^\Delta(F)$ 定义于 6.2（公式 6.2.2）。
\end{SGAresult}
\begin{proof}
证明分为若干步。

a) 设 $\Psi$ 是 $C$ 上的一个\Term{bounded-below}{左有界} $\Lambda[G]$-\Term{module}{模}\Term{sheaf}{层}\Term{complex}{复形}。有\Term{isomorphism}{同构}
\[
\mathrm R\Gamma^G(\mathrm R\Gamma_C(\Psi))
\simeq\mathrm R\Gamma_C(\mathrm R\underline\Gamma^G(\Psi)),\tag{7.3}
\]
这里 $\underline\Gamma^G$ 是“取由 $G$-\Term{invariant}{不变}\Term{section}{截面}组成的\Term{subsheaf}{子层}”的\Term{functor}{函子}。
事实上，$\Gamma_C$ 把\Term{injective-sheaf}{内射} $\Lambda[G]$-\Term{module}{模}\Term{sheaf}{层}变为\Term{injective-module}{内射} $\Lambda[G]$-\Term{module}{模}，而 $\underline\Gamma^G$ 把\Term{injective-sheaf}{内射} $\Lambda[G]$-\Term{module}{模}\Term{sheaf}{层}变为\Term{injective-sheaf}{内射} $\Lambda$-\Term{module}{模}\Term{sheaf}{层}（例如可以利用 $\Gamma_C$ 和 $\underline\Gamma^G$ 的显然\Term{adjoint-functor}{伴随函子}来证明）。另一方面，有
\[
\Gamma^G\circ\Gamma_C=\Gamma_C\circ\underline\Gamma^G.
\]
对两边取\Term{derived-functor}{导出函子}并应用（[5] 命题 3.1），即得（7.3）。

b) 设 $\Phi$ 是 $C$ 上的\Term{bounded-below}{左有界} $\Lambda$-\Term{module}{模}\Term{sheaf}{层}\Term{complex}{复形}，在 $U$ 外为零。则有\Term{isomorphism}{同构}
\[
\Phi\simeq\mathrm R\underline\Gamma^G(p_*(\Phi')),
\qquad\Phi'=p^*(\Phi).\tag{7.4}
\]
先注意，\Term{functor}{函子} $p^*$、$\underline\Gamma^G\circ p_*$ 在下列两个\Term{category}{范畴}间给出彼此\Term{quasi-inverse}{拟逆}的\Term{equivalence}{等价}：$C'$ 上带 $G$ \Term{action}{作用}且在 $U'=p^{-1}(U)$ 外为零的 $\Lambda$-\Term{module}{模}\Term{sheaf}{层}所成的\Term{category}{范畴}，以及 $C$ 上在 $U$ 外为零的 $\Lambda$-\Term{module}{模}\Term{sheaf}{层}所成的\Term{category}{范畴}。所以有\Term{isomorphism}{同构}
\SourcePage{410}
\[
\Phi\xrightarrow{\sim}\mathrm R\bigl((\underline\Gamma^G\circ p_*)\circ p^*\bigr)(\Phi).
\]
另一方面，可以再次应用（[5] 命题 3.1），考虑到 $\mathrm Rp^*=p^*$、$\mathrm Rp_*=p_*$，得到
\[
\Phi\xrightarrow{\sim}\mathrm R(\underline\Gamma^G\circ p_*)(p^*(\Phi))
\simeq(\mathrm R\underline\Gamma^G)(p_*p^*(\Phi)),
\]
这就是所要的\Term{isomorphism}{同构}。

c) 设 $\Phi$ 为 b) 中的\Term{sheaf}{层}\Term{complex}{复形}。有\Term{isomorphism}{同构}
\[
\mathrm R\Gamma_C(\Phi)\simeq
(\mathrm R\Gamma^G)(\mathrm R\Gamma_{C'}(\Phi')),
\qquad\Phi'=p^*(\Phi).\tag{7.5}
\]
事实上，由（7.3）有
\[
\begin{aligned}
(\mathrm R\Gamma^G)(\mathrm R\Gamma_{C'}(\Phi'))
&\simeq(\mathrm R\Gamma^G)(\mathrm R\Gamma_C(p_*(\Phi')))\\
&\simeq\mathrm R\Gamma_C(\mathrm R\underline\Gamma^G(p_*(\Phi'))),
\end{aligned}
\]
再由\Term{isomorphism}{同构}（7.4）得到（7.5）。

d) 用 $i':U'\hookrightarrow C'$、$i:U\hookrightarrow C$ 表示\Term{canonical}{典范}包含，令
\[
\Phi=F_{U,C}=i_!(F|U).
\]
为计算 $\mathrm R\Gamma_C(\Phi)$，可以应用（7.5）。另一方面，按 7.1 的记号，有 $\Phi\simeq\Lambda_{U',C'}\otimes_\Lambda M$，\Term{kunneth-formula}{Künneth 公式}（SGA 4 XVII）给出
\[
\begin{aligned}
\mathrm R\Gamma_{C'}(\Phi')
&\simeq\mathrm R\Gamma_{C'}(\Lambda_{U',C'}\otimes_\Lambda^{\mathrm L}M)\\
&\xleftarrow{\sim}\mathrm R\Gamma_{C'}(\Lambda_{U',C'})\otimes_\Lambda^{\mathrm L}M.
\end{aligned}\tag{7.6}
\]
\SourcePage{411}
最后这个\Term{isomorphism}{同构}实际上是两边所定义的 $\mathrm D(\Lambda[G])$ 的\Term{object}{对象}之间的\Term{isomorphism}{同构}。结合（7.6）、（7.5）与公式（3.9.1）（由 2.1 可知此处可以应用），得到
\[
\begin{aligned}
\mathrm R\Gamma_C(F_{U,C})
&\simeq\mathrm R\Gamma^G\bigl(\mathrm R\Gamma_{C'}(\Lambda_{U',C'})\otimes_\Lambda^{\mathrm L}M\bigr)\\
&\simeq\mathrm R\operatorname{Hom}_{\Lambda[G]}
\bigl(\mathrm R\Gamma_{C'}(\Lambda_{U',C'})^\vee,M\bigr).
\end{aligned}\tag{7.7}
\]
所以，由 $M\simeq F_{\bar\eta}\in\operatorname{Ob}(\Delta)$（条件 i)）及（3.6）得到
\[
\mathrm R\Gamma_C(F_{U,C})\in\operatorname{Ob}(\Delta).
\]
另一方面，有\Term{distinguished-triangle}{区分三角}（其中 $Y=C-U$）
\[
\begin{tikzcd}
&\mathrm R\Gamma_Y(F|Y)\arrow[dl]&\\
\mathrm R\Gamma_C(F_{U,C})\arrow[rr]&&\mathrm R\Gamma_C(F)\arrow[ul].
\end{tikzcd}\tag{7.8}
\]
这里显然有 $\mathrm R\Gamma_Y(F)=\prod_{y\in Y}F_{\bar y}\in\operatorname{Ob}(\Delta)$，因为 $F_{\bar y}\in\operatorname{Ob}(\Delta)$（条件 i)）。因此 $\mathrm R\Gamma_C(F)\in\operatorname{Ob}(\Delta)$，即得\Term{theorem}{定理}的第一个断言。\Term{distinguished-triangle}{区分三角}（7.8）还给出关系
\[
\operatorname{cl}_{K(\Delta)}(\mathrm R\Gamma_C(F))
=\operatorname{cl}_{K(\Delta)}(\mathrm R\Gamma_C(F_{U,C}))
+\sum_{y\in Y}\operatorname{cl}_{K(\Delta)}(F_{\bar y}).
\]
再利用\Term{isomorphism}{同构} 7.7，得到
\[
\begin{aligned}
\operatorname{cl}_{K(\Delta)}(\mathrm R\Gamma_C(F))
={}&\operatorname{cl}_{K(\Delta)}\!\left(\mathrm R\operatorname{Hom}_{\Lambda[G]}
(\mathrm R\Gamma_{C'}(\Lambda_{U',C'})^\vee,F_{\bar\eta})\right)\\
&+\sum_{y\in Y}\operatorname{cl}_{K(\Delta)}(F_{\bar y}).
\end{aligned}
\]
\SourcePage{412}
现在只须使用 Weil 公式（5.2），就得到 \Term{euler-poincare-formula}{Euler--Poincaré 公式}。事实上，首先由 4.4、（5.2）以及 3.6（取 $B=\Lambda[G]$ 的特殊情形），得到
\[
\operatorname{cl}_{K^\bullet(\Lambda[G])}(\mathrm R\Gamma_{C'}(\Lambda_{U',C'}))
=\operatorname{cl}_{K^\bullet(\Lambda[G])}(\mathrm R\Gamma_{C'}(\Lambda_{U',C'})^\vee).
\]
于是
\[
\begin{aligned}
\operatorname{cl}_{K(\Delta)}(\mathrm R\Gamma_C(F))
={}&\operatorname{hom}\Bigl(
\operatorname{cl}_{K^\bullet(\Lambda[G])}(\mathrm R\Gamma_{C'}(\Lambda_{U',C'})),
\operatorname{cl}_{K(\mathrm D(\Lambda[G])_\Delta)}(F_{\bar\eta})\Bigr)
+\sum_{y\in Y}\operatorname{cl}_{K(\Delta)}(F_{\bar y})\\
={}&\sum_{y\in Y}\operatorname{cl}_{K(\Delta)}(F_{\bar y})
+\operatorname{hom}\Bigl(
\operatorname{EP}(C)\operatorname{cl}_{K^\bullet(\Lambda[G])}(\Lambda[G])\\
&\hspace{24mm}-\sum_{y\in Y}
\bigl(\operatorname{sw}_y^\Lambda+\operatorname{cl}_{K^\bullet(\Lambda[G])}(\Lambda[G])\bigr),
\operatorname{cl}_{K(\mathrm D(\Lambda[G])_\Delta)}(F_{\bar\eta})\Bigr)\\
={}&\sum_{y\in Y}\operatorname{cl}_{K(\Delta)}(F_{\bar y})
+\operatorname{EP}(C)\operatorname{cl}_{K(\Delta)}(F_{\bar\eta})
-\sum_{y\in Y}\bigl(\alpha_y^\Delta(F)+\operatorname{cl}_{K(\Delta)}(F_{\bar\eta})\bigr)\\
={}&\operatorname{EP}(C)\operatorname{cl}_{K(\Delta)}(F_{\bar\eta})
-\sum_{y\in Y}\varepsilon_y^\Delta(F).
\end{aligned}
\]
考虑到 $x\in U$ 时 $\varepsilon_x^\Delta(F)=0$，这正是公式（7.2）。
\end{proof}
对于满足 7.1 的条件的\Term{sheaf}{层}\Term{complex}{复形} $F$，引入记号
\[
\operatorname{EP}_\Delta(C,F)=\operatorname{cl}_{K(\Delta)}(\mathrm R\Gamma_C(F)).
\]

\SourcePage{413}
\begin{SGAresult}{推论}{7.9}
若 $F,F'$ 是满足 7.1 的条件的\Term{sheaf}{层}\Term{complex}{复形}，并且它们作为 $A$-\Term{module}{模}\Term{complex}{复形}\Term{local-scheme}{局部}\Term{isomorphism}{同构}（相对于\Term{etale-topology}{平展拓扑}），则
\[
\operatorname{EP}_\Delta(C,F)=\operatorname{EP}_\Delta(C,F').\tag{7.10}
\]
特别地，若 $F$ 是取值 $M^\bullet$ 的\Term{locally-constant}{局部常值}\Term{complex}{复形}，则
\[
\operatorname{EP}_\Delta(C,F)=\operatorname{EP}(C)\operatorname{cl}_{K(\Delta)}(M^\bullet).\tag{7.11}
\]
\end{SGAresult}

\begin{SGAresult}{推论}{7.12}
设 $C$ 是 $k$ 上的\Term{connected}{连通}、\Term{smooth}{光滑}、\Term{proper-scheme}{固有}\Term{algebraic-curve}{代数曲线}，$V$ 是 $C$ 的\Term{nonempty}{非空}\Term{open-subset}{开集}，$\Delta$ 是如 7.1 中那样的 $\mathrm D^+(A)$ 的\Term{subcategory}{子范畴}，$F$ 是 $V$ 上的 $A$-\Term{module}{模}\Term{sheaf}{层}\Term{complex}{复形}，且 $F$ 在 $V$ 的某个适当\Term{nonempty}{非空}子\Term{open-subset}{开集}上的\Term{restriction}{限制}可由 $V$ 的一个\Term{galois-covering}{Galois 覆盖}\Term{trivialization}{平凡化}。又假设对每个 $x\in V$，有 $F_{\bar x}\in\operatorname{Ob}(\Delta)$。则
\[
\mathrm R_!\Gamma_V(F),\ \mathrm R\Gamma_V(F)\in\operatorname{Ob}(\Delta),\tag{*}
\]
\[
\operatorname{cl}_{K(\Delta)}(\mathrm R_!\Gamma_V(F))
=\operatorname{cl}_{K(\Delta)}(\mathrm R\Gamma_V(F)).\tag{**}
\]
令
\[
\operatorname{EP}_\Delta(V,F)=\operatorname{cl}_{K(\Delta)}(\mathrm R\Gamma_V(F)),
\]
以及
\[
\operatorname{EP}(V)=\operatorname{EP}_{\mathrm D^b(\mathbf F_\ell)_{\mathrm{coh}}}(V,\mathbf F_\ell^\bullet)
=\operatorname{EP}(C)-\operatorname{card}(Z),
\]
其中 $Z=C-V$，$\mathbf F_\ell^\bullet$ 是 $C$ 上由与 $\mathbf F_\ell=\mathbf Z/\ell\mathbf Z$ 对应的\Term{constant-sheaf}{常值层}定义的、集中于零次的\Term{complex}{复形}。则有公式
\SourcePage{414}
\[
\operatorname{EP}_\Delta(V,F)
=\operatorname{EP}(V)\operatorname{cl}_{K(\Delta)}(F_{\bar\eta})
-\sum_{x\in V(k)}\varepsilon_x^\Delta(F)
-\sum_{x\in Z}\alpha_x^\Delta(F),\tag{7.13}
\]
其中 $\varepsilon_x^\Delta$ 和 $\alpha_x^\Delta$ 由（6.2.1）、（6.2.2）定义。
\end{SGAresult}

我们将使用下述\Term{lemma}{引理}。
\begin{SGAresult}{\Term{lemma}{引理}}{7.14}
设 $K$ 是\Term{strictly-local}{严格局部}的\Term{discrete-valuation-ring}{离散赋值环} $R$ 的\Term{fraction-field}{分式域}，$W=\operatorname{Spec}K$。设 $A$ 是被整数 $n$ \Term{annihilate}{零化}的\Term{noetherian}{诺特}环，$n$ 与 $R$ 的\Term{residue-characteristic}{剩余特征} $p$ 互素；$F$ 是 $W$ 上的 $A$-\Term{module}{模}\Term{sheaf}{层}\Term{complex}{复形}，可由 $K$ 的一个\Term{finite-extension}{有限扩张}\Term{trivialization}{平凡化}。若 $\Delta$ 为 $\mathrm D^+(A)$ 的一个对\Term{direct-summand}{直和因子}\Term{stable}{稳定}的\Term{triangulated-subcategory}{三角子范畴}，且 $F_{\bar\eta}\in\operatorname{Ob}(\Delta)$（$\eta$ 为 $S$ 的\Term{generic-point}{泛点}），则
\[
\mathrm R\Gamma_W(F)\in\operatorname{Ob}(\Delta),\qquad
\operatorname{cl}_{K(\Delta)}(\mathrm R\Gamma_W(F))=0.
\]
\end{SGAresult}
\begin{proof}[7.14 的证明]
设 $\pi$ 为 $W$ 的\Term{fundamental-group}{基本群}。则 $F$ 对应一个以 $\pi$ 为\Term{operator-group}{算子群}的 $A$-\Term{module}{模}\Term{complex}{复形} $M^\bullet$。$M^\bullet$ 的底层 $A$-\Term{module}{模}\Term{complex}{复形}与 $F_{\bar\eta}$ \Term{isomorphism}{同构}，所以属于 $\operatorname{Ob}(\Delta)$。引入 $\pi$ 对其唯一的 Sylow $p$-\Term{subgroup}{子群} $P$ 的\Term{quotient-group}{商群} $\bar\pi$ [4，II 33]，也就是 $\pi$ 的所有阶与 $p$ 互素的\Term{discrete}{离散}\Term{quotient-group}{商群}的\Term{projective-limit}{投射极限}，这里 $p$ 是 $R$ 的\Term{residue-field}{剩余域}的\Term{characteristic-exponent}{特征指数}。（若 $p=1$，则置 $\bar\pi=\pi$。）所以有\Term{exact-sequence}{正合列}
\[
1\longrightarrow P\longrightarrow\pi\longrightarrow\bar\pi\longrightarrow0,
\]
并得到
\[
\mathrm R\Gamma_W(F)\simeq\mathrm R\Gamma^\pi(M^\bullet)
\simeq\mathrm R\Gamma^{\bar\pi}(\mathrm R\Gamma^P(M^\bullet))
\]
\SourcePage{415}
（“Hochschild--Serre \Term{isomorphism}{同构}”）。但 $P$ 通过一个\Term{finite}{有限}\Term{quotient-group}{商群}\Term{action}{作用}于 $M^\bullet$；该\Term{quotient-group}{商群}是 $p$-群，故其阶与 $n$ 互素。因此
$\mathrm R\Gamma^P(M^\bullet)\simeq\Gamma^P(M^\bullet)$，而 $\Gamma^P(M^\bullet)$ 为 $M^\bullet$ 的\Term{direct-summand}{直和因子}，所以属于 $\operatorname{Ob}(\Delta)$。

另一方面，$\bar\pi=\prod_{\ell\ne p}\mathbf Z_\ell$（[4]，II 33）。因此归结为证明：若 $M^\bullet\in\operatorname{Ob}(\Delta)$，则 $\mathrm R\Gamma^{\bar\pi}(M^\bullet)\in\operatorname{Ob}(\Delta)$，且 $\operatorname{cl}_{K(\Delta)}(\mathrm R\Gamma^{\bar\pi}(M^\bullet))=0$。为此，取 $\bar\pi$ 的一个\Term{topological-generator}{拓扑生成元} $f$。可以假设 $M$ 的各\Term{set-component}{分量} $M^i$ 是满足 $j\geq1$ 时 $H^j(\bar\pi,M^i)=0$ 的 $\bar\pi$-\Term{module}{模}。这时有\Term{distinguished-triangle}{区分三角}
\[
\begin{tikzcd}
&M^\bullet\arrow[dl]&\\
\mathrm R\Gamma^{\bar\pi}(M^\bullet)\arrow[rr]&&M^\bullet\arrow[ul,"1-f"'].
\end{tikzcd}\tag{7.15}
\]
事实上，一方面显然 $\mathrm R\Gamma^{\bar\pi}=\operatorname{Ker}(1-f)$；另一方面，$1-f:M^\bullet\longrightarrow M^\bullet$ 是\Term{surjectivity}{满射}，因为
\[
H^1(\bar\pi,M)=M^\bullet/(1-f)M^\bullet
\]
（这一计算可参看 [4] 第 197 页中一个大致相似的情形），所以 $M^\bullet=(1-f)M^\bullet$。\Term{distinguished-triangle}{三角}（7.15）同时给出
\[
\mathrm R\Gamma^{\bar\pi}(M^\bullet)\in\operatorname{Ob}(\Delta),\qquad
\operatorname{cl}_{K(\Delta)}(\mathrm R\Gamma^{\bar\pi}(M^\bullet))=0,
\]
从而证明了 7.14。
\end{proof}

\SourcePage{416}
\begin{proof}[7.12 的证明]
对\Term{complex}{复形} $i_!(F)$ 应用 7.1，立即得到 $\mathrm R_!\Gamma_V(F)\in\operatorname{Ob}(\Delta)$；这里 $i:V\hookrightarrow C$，且
$\mathrm R_!\Gamma_V(F)=\mathrm R\Gamma_C(i_!(F))$。为证明 $\mathrm R\Gamma_V(F)\in\operatorname{Ob}(\Delta)$，考虑\Term{distinguished-triangle}{区分三角}
\[
\begin{tikzcd}
&Q^\bullet\arrow[dl]&\\
\mathrm Ri_!(F)\arrow[rr]&&\mathrm Ri_*(F)\arrow[ul],
\end{tikzcd}
\]
其中 $Q^\bullet$ 集中于 $Z=C-V$ 上。对它应用 $\mathrm R\Gamma_C$，得\Term{distinguished-triangle}{区分三角}
\[
\begin{tikzcd}
&\mathrm R\Gamma_Z(Q^\bullet|Z)\arrow[dl]&\\
\mathrm R_!\Gamma_V(F)\arrow[rr]&&\mathrm R\Gamma_V(F)\arrow[ul].
\end{tikzcd}\tag{7.16}
\]
由于 $\mathrm R_!\Gamma_V(F)\in\operatorname{Ob}(\Delta)$，可知 $\mathrm R\Gamma_V(F)\in\operatorname{Ob}(\Delta)$ 当且仅当 $\mathrm R\Gamma_Z(Q^\bullet|Z)\in\operatorname{Ob}(\Delta)$，也就是对每个 $x\in Z$ 都有 $(\mathrm Ri_*(F))_{\bar x}\in\operatorname{Ob}(\Delta)$。回顾一下，计算 $(\mathrm Ri_*(F))_{\bar x}$ 时，要考虑\Term{cartesian-square}{笛卡尔方块}
\[
\begin{tikzcd}
\overline V\arrow[r,"j"]\arrow[d,"u"']&C_x\arrow[d]\\
V\arrow[r]&C,
\end{tikzcd}
\]
其中 $C_x$ 为对应于\Term{geometric-point}{几何点} $x$ 的\Term{strict-localization}{严格局部化}\Term{scheme}{概形}。有
\[
(\mathrm Ri_*(F))_{\bar x}\simeq\mathrm R\Gamma_{\overline V}(u^*(F)).
\]
由 7.14，
\SourcePage{417}
于是有
\[
(\mathrm Ri_*(F))_{\bar x}\in\operatorname{Ob}(\Delta),\qquad
\operatorname{cl}_{K(\Delta)}((\mathrm Ri_*(F))_{\bar x})=0.
\]
所以 $\operatorname{cl}_{K(\Delta)}(\mathrm R\Gamma_Z(Q^\bullet|Z))=0$，而\Term{distinguished-triangle}{区分三角}（7.16）给出
\[
\operatorname{cl}_{K(\Delta)}(\mathrm R_!\Gamma_V(F))
=\operatorname{cl}_{K(\Delta)}(\mathrm R\Gamma_V(F)).
\]
最后，要证明（7.13），只须对\Term{sheaf}{层}\Term{complex}{复形} $i_!(F)$ 应用 \Term{euler-poincare-formula}{Euler--Poincaré 公式}（7.2）：
\[
\begin{aligned}
\operatorname{EP}_\Delta(V,F)
&=\operatorname{cl}_{K(\Delta)}(\mathrm R_!\Gamma_V(F))
=\operatorname{cl}_{K(\Delta)}(\mathrm R\Gamma_C(i_!(F)))\\
&=\operatorname{EP}(C)\operatorname{cl}_{K(\Delta)}(F_{\bar\eta})
-\sum_{x\in C(k)}\varepsilon_x^\Delta(i_!(F)).
\end{aligned}
\]
而
\[
\sum_{x\in C(k)}\varepsilon_x^\Delta(i_!(F))
=\sum_{x\in V(k)}\varepsilon_x^\Delta(F)
+\sum_{x\in Z}\varepsilon_x^\Delta(F).
\]
但 $x\in Z$ 时 $(i_!(F))_{\bar x}=0$，因而对于这样的 $x$，有
\[
\varepsilon_x^\Delta(i_!(F))
=\alpha_x^\Delta(F)+\operatorname{cl}_{K(\Delta)}(F_{\bar\eta}),
\]
由此得到公式（7.13）。
\end{proof}

\SourcePage{418}
\SGAsection{}{参考文献}
\begin{description}
\item[{[1]}] M. Raynaud，\emph{层的 \Term{euler-poincare-characteristic}{Euler--Poincaré 特征数}与\Term{abelian-variety}{阿贝尔簇}的\Term{cohomology}{上同调}}（Caractéristique d'Euler--Poincaré d'un faisceau et cohomologie des variétés abéliennes），Séminaire Bourbaki，1965，报告第 286 号。
\item[{[2]}] A. Grothendieck，\emph{\Term{homological-algebra}{同调代数}的若干问题}（Sur quelques points d'algèbre homologique），Toh. Math. Journal，第 9 卷（1957），第 119--221 页（文中引作 [T]）。
\item[{[3]}] J.-P. Serre，\emph{\Term{local-scheme}{局部}域}（Corps locaux），Act. Sci. Ind. 1296（1962），Hermann，Paris。
\item[{[4]}] J.-P. Serre，\emph{Galois \Term{cohomology}{上同调}}（Cohomologie Galoisienne），Lecture Notes in Mathematics 5，1964（Springer）。
\item[{[5]}] J.-L. Verdier，\emph{\Term{derived-category}{导出范畴}}（Catégories dérivées），IHES，1963（发表于 SGA $4\frac12$）。
\end{description}
